\section{What is shared, and what changes in Lean}\label{sec:comparison}
The common proof does not make the two native theorem statements identical.
Table~\ref{tab:comparison} records the relevant distinctions.
\begin{table}[ht]
\caption{The two verified instantiations. The weights refer to physical coefficients.}\label{tab:comparison}
\centering
\begin{tabular}{@{}p{0.27\textwidth}p{0.31\textwidth}p{0.31\textwidth}@{}}
\toprule
Component & Two dimensions & Four dimensions\\
\midrule
Physical site covariance & $a^{-2}G_N$ & $a^{-4}G_N$\\
Raw lift & $a^2\sum_x\phi_x\delta_{ax}$ & $a^4\sum_x\phi_x\delta_{ax}$\\
Moment / limit spaces & $H^{-2}$ / $H^{-3}$ & $H^{-3}$ / $H^{-4}$\\
Native limit weight & $1$ at zero; $R(k)^{-6}$ otherwise & $(1+R(k))^{-8}$\\
Limit construction & Prokhorov existence plus identified cylinders & Named explicit Gaussian-series law\\
Identification & All finite weighted complex cylinders & Canonical real cylinders and real support\\
Sequence & $N=n+1$, $a=L/N$ & $N=n+1$, $a=L/N$\\
\bottomrule
\end{tabular}
\end{table}

\subsection{The native two-dimensional theorem}
The final two-dimensional result asserts existence of a Borel probability
measure on the concrete completion, convergence of the full lattice sequence
to it, and equality of each finite Fourier cylinder with the corresponding
periodic-resolvent Gaussian marginal. In the source, the main declaration is
\begin{quote}\raggedright
\code{periodicHMinusThree_fixedVolume_weak_limit_exists_with_expected_cylinders}.
\end{quote}
The finite cylinder reads the stored weighted complex coordinates. Its target
is consequently obtained by complexifying cosine/sine Gaussian evaluations
and multiplying each by the square root of the native weight. Dividing by
that nonzero square root recovers physical Fourier marginals. The theorem
has no caller-supplied measurability or tightness input.

Here the native stronger and weaker weights are
\[
 v_r(k)=\begin{cases}1,&k=0,\\R(k)^{-2r},&k\ne0,\end{cases}
 \qquad r=2,3.
\]
For integer frequencies, $w_r\leq v_r\leq2^{2r}w_r$. Thus the native carrier
is topologically equivalent to our common convention, but the norms are
not literally equal. The full $H^{-2}$ energy includes the zero mode; a
zero-mode-free seminorm would not control the stated Gaussian family.

\subsection{The native four-dimensional theorem}
The four-dimensional statement has the more explicit form
\begin{quote}\ttfamily\raggedright
fs4PeriodicFixedVolume\_weak\_tendsto\_free.
\end{quote}
It proves convergence to the named law of a measurable Gaussian series on
the full weighted completion. Positivity of $L$ and $m$ are its only explicit
model hypotheses. The proof identifies every finite injective family of
canonical real coordinates with actual periodic-resolvent Gaussian targets,
passes real support to weak sublimits, and invokes real-supported cylinder
uniqueness. The arbitrary-complex-carrier extensionality claim would be
false for that reduced coordinate set and is not used.

The rank-four summability argument splits the integer lattice into two
rank-two blocks. Mathlib's norm on a finite function type is a supremum norm:
$\norm{(1,1)}_\infty=1$, not $\sqrt2$. The block expression in the weight
therefore equals the coordinate maximum $R(k)$. This is compatible with a
Euclidean covariance symbol because equivalent carrier norms do not change
the topology. They must not be confused with equality of norms or rotation
isometries.

\subsection{Reuse at the level of mathematical contracts}
Both developments use the same concrete finite-lattice covariance layer,
physical test evaluations and periodic covariance limit. They then prove their
own carrier membership, stronger-energy bound and compactness instance.
The reusable contract is: construct actual probability laws on a Polish
carrier, establish a stronger compactly embedded moment bound, identify all
finite Gaussian evaluations, and prove those evaluations determine the law.
The article abstracts this contract in Theorem~\ref{thm:main}; the current
Lean library proves its two specified realizations. A generic Lean theorem
for arbitrary dimension would require additional parametrized carrier and
summability work.
